\documentclass{article}
\usepackage[utf8]{inputenc}
\usepackage{amsmath}
\usepackage{amsfonts}
\usepackage{amssymb}

\title{Documento Exemplo com Seções e Subseções}
\author{Assistente de IA}
\date{\today}

\begin{document}

\maketitle

\tableofcontents
\newpage

\section{Introdução}
Este é o conteúdo da seção de introdução. Aqui apresentamos o propósito e o escopo do documento.

\subsection{Objetivo do Documento}
Detalhes sobre o objetivo principal deste documento.

\subsection{Estrutura do Documento}
Explicação sobre como o documento está organizado.

\section{Fundamentação Teórica}
Nesta seção, abordaremos os conceitos teóricos fundamentais.

\subsection{Conceitos Básicos}
Definições e conceitos iniciais importantes.

\subsubsection{Definições Preliminares}
Detalhamento das definições básicas.

\subsubsection{Notação Utilizada}
Explicação da notação que será utilizada ao longo do documento.

\subsection{Teorias Relevantes}
Apresentação das teorias que fundamentam o trabalho.

\section{Metodologia}
Descrição da metodologia utilizada no desenvolvimento do trabalho.

\subsection{Etapas de Desenvolvimento}
Detalhamento das etapas seguidas.

\subsection{Ferramentas Utilizadas}
Lista e descrição das ferramentas utilizadas.

\section{Resultados}
Apresentação e análise dos resultados obtidos.

\subsection{Análise Quantitativa}
Resultados numéricos e estatísticos.

\subsection{Análise Qualitativa}
Resultados descritivos e observacionais.

\section{Conclusão}
Resumo dos principais pontos abordados e considerações finais.

\subsection{Contribuições}
Principais contribuições deste trabalho.

\subsection{Trabalhos Futuros}
Sugestões para trabalhos futuros baseados neste.

\end{document}